Project identity (name:Grzegorz) 

 

Logo and Branding Design (Ola) 

The \abbr{EWAS} mission logo serves as a straightforward visual identifier for the project, using established aerospace symbols and functional colors to reflect the system's civil protection role. 

 

Shape and Geometry 

    The Rounded Triangle Frame: Chosen to mimic standard warning and alert signs, immediately communicating the early warning nature of the architecture. 

    The Spacecraft \& Earth Backdrop: The central illustration depicts a satellite with deployed solar arrays, while the blue background symbolizes the operational environment. 

    The Conical Beam: A yellow cone projecting from the payload visualizes the continuous footprint of the Short-Wave (\abbr{SWIR}) and Medium-Wave (\abbr{MWIR}) infrared sensors monitoring the protected region. 

Color Palette Meaning 

    Deep Blue \& Light Blue: Serving as the primary visual theme of the logo, these blue tones represent space surveillance, high-altitude operations, and the Earth's atmosphere. More importantly, blue is intentionally used to communicate reliability and trust—two critical psychological factors for an infrastructure designed for national crisis management. 

    Safety Yellow: Used exclusively for the sensor beam to represent active detection, tracking, and the conversion of thermal data into a live digital alert. Additionally, the color yellow reinforces the overall warning and alert nature of the project. 

    Metallic Gray / Silver: Applied to the satellite bus to signify the hardware-driven nature of the project and the use of standardized, flight-proven components. Furthermore, gray symbolizes technology and modernity, emphasizing how revolutionary this space-based early warning solution is compared to traditional radar networks. 

    Black: Used for the heavy border and the bold "\abbr{EWAS}" typography to give the logo a cleaner and bolder look. Strategically, this strong presence of black signals the operational power and tactical advantage gained by national decision-makers through the highly useful, low-latency information delivered by the system. 